\documentclass[11pt]{article}
\usepackage[margin=1in]{geometry}
\usepackage{amsmath,amssymb,amsthm,mathtools}
\usepackage{array,booktabs}
\usepackage{graphicx}
\usepackage{xcolor}
\usepackage[hidelinks]{hyperref}
\usepackage{microtype}

\definecolor{accent}{RGB}{37,86,123}
\newtheorem{theorem}{Theorem}
\newtheorem{lemma}[theorem]{Lemma}
\newtheorem{corollary}[theorem]{Corollary}
\theoremstyle{remark}
\newtheorem{remark}[theorem]{Remark}

\newcommand{\C}{\mathbb C}
\newcommand{\cB}{\mathcal B}
\newcommand{\cb}{\mathrm{cb}}
\newcommand{\ran}{\operatorname{ran}}

\title{\textcolor{accent}{Equality of the two continuous
Fourier--Stieltjes norms for discrete groupoids}}
\author{Automated mathematics run \texttt{fa\_banach\_001}, lane 13}
\date{23 July 2026}

\begin{document}
\maketitle

\begin{center}
\fbox{\parbox{0.88\textwidth}{
\textbf{Status.} Complete theorem for the discrete subclass of the source
question; current verdict: likely valid.  The nondiscrete problem remains open.
}}
\end{center}

\section{Source question and result}

Let \(G\) be a locally compact groupoid with a Haar system.  Paterson
\cite[\S5]{Paterson2003} defines the continuous Fourier--Stieltjes algebra
\(B(G)\) from coefficients of continuous \(G\)-Hilbert bundles.  With the
convention that Hilbert-space inner products are conjugate-linear in the first
variable, a coefficient has the form
\[
 \phi(g)=\bigl(L_g\xi(s(g)),\eta(r(g))\bigr),
\]
where \(\xi,\eta\) are bounded continuous sections, and
\[
 \|\phi\|_B
 =\inf\{\|\xi\|_\infty\|\eta\|_\infty:
             \phi=(\xi,\eta)\}.
\]
Pointwise multiplication by \(\phi\) extends to a completely bounded map
\[
 T_\phi:C^*(G)\longrightarrow C^*(G),\qquad
 T_\phi(f)=\phi f,
\]
and one sets \(\|\phi\|_{\cb}=\|T_\phi\|_{\cb}\).  Paterson records the
always-valid inequality
\begin{equation}\label{eq:known}
 \|\phi\|_{\cb}\leq\|\phi\|_B
\end{equation}
and asks whether equality holds in general.  The source gives equality for the
trivial pair groupoid \(R_n\), as well as a separate proof for locally compact
abelian groups.

The exact source sentence (Section~5, page~10 of the paper) is:
\begin{quote}
\emph{``It is an open question if the two norms actually coincide in
general.''}
\end{quote}
The surrounding source passage is reproduced below.

\begin{center}
\includegraphics[width=0.94\textwidth]{figures/open_problem_crop.png}
\end{center}

The theorem below supplies the reverse inequality throughout the discrete
case.  As usual in the source setting, ``discrete'' plus second countability
means countable; equivalently, one can take countability as an explicit
hypothesis.

\section{Proof intuition}

The cb multiplier itself contains the desired coefficient factorization once
it is observed through the trivial one-dimensional groupoid representation.
Indeed, that representation sends an arrow \(g:v\to u\) to the matrix unit
\(E_{u,v}\).  A Haagerup--Paulsen factorization of the resulting
operator-valued map can then be cut by the mutually orthogonal unit
projections of \(C^*(G)\).  Those cuts are the Hilbert fibers, the images of
the arrow point masses are the fiberwise unitaries, and the columns of the two
factorizing operators are bounded coefficient sections.  This reconstructs
\(\phi\) at cost at most \(\|T_\phi\|_{\cb}\), precisely reversing
\eqref{eq:known}.  Discreteness is what makes the point masses and the
pointwise fiber columns available; the nondiscrete source question remains
open.

\begin{theorem}\label{thm:main}
Let \(G\) be a countable discrete groupoid, equipped with its counting Haar
system.  Then
\[
 \|\phi\|_B=\|T_\phi\|_{\cb}
 \qquad\text{for every }\phi\in B(G).
\]
\end{theorem}

\section{The trivial line-bundle representation}

Write \(G^0\) for the unit space and put
\[
 A=C^*(G),\qquad H_0=\ell^2(G^0).
\]
For \(u\in G^0\), let \(e_u\) be the corresponding unit vector.  For
\(g\in G\), write \(\delta_g\in C_c(G)\) for the point mass and
\(p_u=\delta_u\).

\begin{lemma}\label{lem:sigma}
There is a \(^{*}\)-representation
\(\sigma:A\to B(H_0)\) satisfying
\[
 \sigma(\delta_g)e_{s(g)}=e_{r(g)},\qquad
 \sigma(\delta_g)e_v=0\quad(v\ne s(g)).
\]
Equivalently, \(\sigma(\delta_g)=E_{r(g),s(g)}\).
\end{lemma}

\begin{proof}
Give every unit the one-dimensional Hilbert fiber \(\C\), and let every arrow
act by the identity map between its source and range fibers.  This is the
trivial \(G\)-Hilbert bundle.  Its integrated representation on
\(\ell^2(G^0)\) has the displayed formula.

One can also check the formula directly.  If \(g,h\) are composable, then
\[
 E_{r(g),s(g)}E_{r(h),s(h)}
 =E_{r(gh),s(gh)};
\]
if they are not composable, the product is zero, and adjoints correspond to
inverses.  Thus the assignment is a \(^{*}\)-representation of \(C_c(G)\).
For \(f\in C_c(G)\), its matrix row and column sums are bounded by the two
terms in the groupoid \(I\)-norm, so the assignment is \(I\)-norm bounded.
It therefore extends to the full groupoid \(C^*\)-algebra.
\end{proof}

The fact that distinct arrows with the same source and range are sent to the
same matrix unit causes no loss: the domain elements \(\delta_g\) remain
distinct, and we will evaluate the resulting operator-valued map separately
on every one of them.

\section{Reconstructing a coefficient from a cb factorization}

We use the standard Haagerup--Paulsen factorization theorem
\cite[Ch.~8]{Paulsen}: if \(S:A\to B(H_0)\) is completely bounded, then for
each \(\varepsilon>0\) there are a Hilbert space \(K\), a
\(^{*}\)-representation \(\pi:A\to B(K)\), and operators
\(V,W:H_0\to K\) such that
\begin{equation}\label{eq:HP}
 S(a)=V^*\pi(a)W,\qquad
 \|V\|\,\|W\|\leq\|S\|_{\cb}+\varepsilon.
\end{equation}

\begin{lemma}\label{lem:reverse}
For a countable discrete groupoid \(G\) and every \(\phi\in B(G)\),
\[
 \|\phi\|_B\leq\|T_\phi\|_{\cb}.
\]
\end{lemma}

\begin{proof}
Fix \(\varepsilon>0\) and compose the multiplier with the representation from
Lemma~\ref{lem:sigma}:
\[
 S=\sigma\circ T_\phi:A\longrightarrow B(H_0).
\]
Then \(\|S\|_{\cb}\leq\|T_\phi\|_{\cb}\).  Choose a factorization
\eqref{eq:HP}.

The mutually orthogonal projections \(\pi(p_u)\) define Hilbert fibers
\[
 K_u=\pi(p_u)K\qquad(u\in G^0).
\]
For every arrow \(g\), the restriction
\[
 U_g=\pi(\delta_g)|_{K_{s(g)}}:K_{s(g)}\longrightarrow K_{r(g)}
\]
is unitary, because
\[
 \delta_g^*\delta_g=p_{s(g)},\qquad
 \delta_g\delta_g^*=p_{r(g)}.
\]
The convolution relations give \(U_gU_h=U_{gh}\) for composable arrows.
Thus \((K_u,U_g)\) is a \(G\)-Hilbert bundle.  Since \(G^0\) is discrete,
there is no continuity obstruction.

Define bounded sections
\[
 \xi(u)=\pi(p_u)We_u,\qquad
 \eta(u)=\pi(p_u)Ve_u.
\]
They satisfy
\begin{equation}\label{eq:sectionbound}
 \|\xi\|_\infty\leq\|W\|,\qquad
 \|\eta\|_\infty\leq\|V\|.
\end{equation}
For \(g\in G\), using Paterson's convention that the inner product is
conjugate-linear in the first variable, we obtain
\begin{align*}
 \bigl(U_g\xi(s(g)),\eta(r(g))\bigr)
 &=\bigl(\pi(\delta_g)We_{s(g)},Ve_{r(g)}\bigr)\\
 &=\bigl(V^*\pi(\delta_g)We_{s(g)},e_{r(g)}\bigr)\\
 &=\bigl(S(\delta_g)e_{s(g)},e_{r(g)}\bigr)\\
 &=\bigl(\phi(g)e_{r(g)},e_{r(g)}\bigr)
 =\overline{\phi(g)}.
\end{align*}
Hence \(\overline\phi\) is a coefficient with the bounds
\eqref{eq:sectionbound}.  Passing to the conjugate Hilbert bundle turns this
coefficient into \(\phi\): if
\((\overline{x},\overline{y})_{\overline K}=(y,x)_K\), then
\[
 \bigl(\overline{U_g}\,\overline{\xi(s(g))},
       \overline{\eta(r(g))}\bigr)_{\overline K}
 =
 \overline{
 \bigl(U_g\xi(s(g)),\eta(r(g))\bigr)_K}
 =\phi(g).
\]
Conjugation preserves all fiber norms.  Consequently,
\[
 \|\phi\|_B
 \leq\|\xi\|_\infty\|\eta\|_\infty
 \leq\|V\|\,\|W\|
 \leq\|S\|_{\cb}+\varepsilon
 \leq\|T_\phi\|_{\cb}+\varepsilon.
\]
Letting \(\varepsilon\downarrow0\) proves the lemma.
\end{proof}

\begin{proof}[Proof of Theorem~\ref{thm:main}]
Combine Lemma~\ref{lem:reverse} with the established inequality
\eqref{eq:known}.
\end{proof}

\begin{corollary}\label{cor:finite}
The two norms coincide for every finite groupoid.
\end{corollary}

\begin{remark}[What this adds to the recorded finite case]
A transitive finite groupoid is isomorphic, after choices of connecting
arrows, to
\[
 I\times H\times I,
\]
where \(H\) is an isotropy group; its full \(C^*\)-algebra is
\(M_{|I|}(C^*(H))\).  Paterson's \(R_n\) is the case \(H=\{e\}\).
Corollary~\ref{cor:finite} allows arbitrary finite \(H\), including
nonabelian isotropy, and arbitrary finite disjoint unions of such components.
\end{remark}

\section{Scope audit and novelty check}

\subsection*{Why discreteness enters}

The proof uses both the atomic vectors \(e_u\in\ell^2(G^0)\) and the isolated
arrow functions \(\delta_g\in C_c(G)\).  They let a single completely bounded
map recover the value \(\phi(g)\) arrow by arrow.  For a nondiscrete unit
space, counting measure need not be Radon and the constructed fiber vectors
need not form continuous sections.  Thus the proof does not establish the
general locally compact statement.

\subsection*{Adversarial checks}

\begin{center}
\begin{tabular}{@{}>{\raggedright\arraybackslash}p{0.29\textwidth}
>{\raggedright\arraybackslash}p{0.62\textwidth}@{}}
\toprule
Potential failure & Check \\
\midrule
Parallel arrows collapse under \(\sigma\) &
Only the output matrix unit repeats; \(S\) is evaluated on each distinct
domain basis element \(\delta_g\), so its scalar is still \(\phi(g)\).\\
The factorizing \(\pi\) is not a bundle representation &
The exact relations for \(p_u,\delta_g\) cut \(K\) into fibers and make every
\(\pi(\delta_g)\) a unitary between the correct fibers.\\
Sections are not bounded or continuous &
Projection gives the uniform estimates \eqref{eq:sectionbound}; every section
over a discrete base is continuous.\\
Inner-product convention conjugates the answer &
The first coefficient is \(\overline\phi\); the conjugate Hilbert bundle gives
\(\phi\) with identical norms.\\
Factorization estimates the wrong norm &
\(\sigma\) is a \(^{*}\)-representation, so
\(\|\sigma T_\phi\|_{\cb}\leq\|T_\phi\|_{\cb}\).\\
\bottomrule
\end{tabular}
\end{center}

\subsection*{Bounded literature search}

On 23 July 2026, searches were run for the exact source question and for
combinations of \emph{discrete groupoid}, \emph{r-discrete groupoid},
\emph{Fourier--Stieltjes}, \emph{coefficient norm}, and
\emph{completely bounded multiplier}.  Oty's earlier work \cite{Oty} studies
the Borel/r-discrete Fourier--Stieltjes algebra with a norm defined through
multipliers on the Borel completion \(M^*(G)\); this is not a statement of
the equality proved here for Paterson's continuous Hilbert-bundle coefficient
norm on \(B(G)\).  Later operator-space duality \cite{Paterson2011} is in the
measured setting, while recent restriction-map work \cite{DG2025} addresses a
different problem.  No exact prior statement was located.  This is a bounded
negative search, so novelty confidence is moderate rather than definitive.

\subsection*{Human-review recommendation}

Send to a specialist reviewer.  The two points deserving the closest check
are the application of Haagerup--Paulsen factorization to
\(\sigma T_\phi\) and the conjugate-bundle calculation.  If accepted, the
theorem is suitable for a concise proposition or short note.

\begin{thebibliography}{9}

\bibitem{Paterson2003}
A. L. T. Paterson,
\emph{The Fourier--Stieltjes and Fourier algebras for locally compact
groupoids},
\href{https://arxiv.org/abs/math/0310138}{arXiv:math/0310138} (2003).

\bibitem{Paulsen}
V. I. Paulsen,
\emph{Completely Bounded Maps and Operator Algebras},
Cambridge Studies in Advanced Mathematics 78,
Cambridge University Press, 2002.

\bibitem{Oty}
K. J. Oty,
\emph{Fourier--Stieltjes algebras of \(r\)-discrete groupoids},
J. Operator Theory \textbf{41} (1999), 175--197.

\bibitem{Paterson2011}
A. L. T. Paterson,
\emph{An operator space duality theorem for the Fourier--Stieltjes algebra of
a locally compact groupoid},
\href{https://arxiv.org/abs/1102.0304}{arXiv:1102.0304} (2011).

\bibitem{DG2025}
J. DeGaetani and M. Ghandehari,
\emph{On the restriction maps of the Fourier and Fourier--Stieltjes algebras
over locally compact groupoids},
Int. Math. Res. Not. IMRN (2025), rnaf057.

\end{thebibliography}

\end{document}
